%% vol3b_ch6_10.tex — Chapters 6–10: Frame, Dialectic, Zombie, Aspect, Rule
%% Part II of Volume III-B: Paradoxes of Mind and Meaning
%% Uses custom commands: \rs, \emerge, \compose, \void, \IS

%%%%%%%%%%%%%%%%%%%%%%%%%%%%%%%%%%%%%%%%%%%%%%%%%%%%%%%%%%%%%%%%%%%%%%%%
%%  CHAPTER 6 — THE FRAME PROBLEM: BACKGROUND IS INFINITE
%%%%%%%%%%%%%%%%%%%%%%%%%%%%%%%%%%%%%%%%%%%%%%%%%%%%%%%%%%%%%%%%%%%%%%%%

\chapter{The Frame Problem --- Background Is Infinite}
\label{ch:frame-problem}

\epigraph{``The frame problem is not just a technical problem in AI.
It is the deepest question of epistemology: how does a finite mind
manage to think at all, given that the relevant background is
infinite?''}{Daniel Dennett, \textit{Cognitive Wheels} (1984)}

\section{The Frame Problem in Philosophy and Artificial Intelligence}

In 1969, John McCarthy and Patrick Hayes identified a puzzle that would
haunt artificial intelligence for the next half-century.  They called it
the \emph{frame problem}: when an agent performs an action, how does it
determine what has \emph{not} changed?  If a robot paints a wall red,
does the wall's weight change?  Does the room's temperature change?
Does the robot's owner's birthday change?  Every action has a finite
number of effects, but the number of things that remain unchanged is
infinite.  How do you specify all of them?

The original frame problem was technical---a matter of efficient
representation in logical AI systems.  But as Dennett argued in his
celebrated essay ``Cognitive Wheels'' (1984), the deeper issue is
philosophical.  Every act of understanding occurs against a
\emph{background} of assumptions so vast that no finite list could
exhaust it.  Hubert Dreyfus, drawing on Heidegger, made a similar
point: the background of ``ready-to-hand'' practices that make
understanding possible is not a set of propositions at all but a
holistic, lived competence that resists articulation.

In this chapter, we prove that the frame problem is not merely a
practical difficulty but a \emph{mathematical inevitability}.  In the
ideatic space, the background generated by a single idea grows
monotonically under self-composition.  You cannot compress it, you
cannot reduce it, and you cannot escape it.  The frame is always
growing, always accumulating, and always constraining its own growth
in ways that are precisely quantifiable.

\section{Self-Composition and the Growth of Background}

We begin with a recursive definition that captures the accumulation
of background through repeated self-composition.

\begin{definition}[Iterated Self-Composition]
\label{def:composeN}
For any idea $a \in \IS$ and natural number $n \in \N$, the
\textbf{$n$-fold self-composition} $\mathrm{composeN}(a, n)$ is defined
recursively:
\begin{align}
\mathrm{composeN}(a, 0) &= \void, \\
\mathrm{composeN}(a, n+1) &= \mathrm{composeN}(a, n) \compose a.
\end{align}
\end{definition}

This definition captures a natural process: the accumulation of
context through repetition.  When you read a word once, you have
$\mathrm{composeN}(\text{word}, 1) = \void \compose \text{word} = \text{word}$.
When you read it again in a new context, you have $\mathrm{composeN}(\text{word}, 2)
= \text{word} \compose \text{word}$---the word composed with itself, which by
Axiom~E4 is \emph{at least} as heavy as the word alone.  Each repetition
adds another layer of background, another fold of context.

The base case $\mathrm{composeN}(a, 0) = \void$ says that zero encounters
leave you with nothing---the empty background.  The recursive case says
that each new encounter composes the current background with the idea
itself.

\begin{lstlisting}
-- From IDT_v3b_DeepPhilosophy.lean

/-- n-fold self-composition: composing idea a with itself n times. -/
def composeN (a : I) : ℕ → I
  | 0 => void
  | n + 1 => compose (composeN a n) a
\end{lstlisting}

\section{The Monotonicity of Background (DP22)}

\begin{theorem}[Background Weight Grows Monotonically --- DP22]
\label{thm:background-weight-monotone}
For any idea $a \in \IS$ and any $n \in \N$:
\[
\mathrm{weight}(\mathrm{composeN}(a, n+1)) \geq
\mathrm{weight}(\mathrm{composeN}(a, n)) \geq 0.
\]
The weight of the accumulated background is non-decreasing and always
non-negative.
\end{theorem}

\begin{proof}
The first inequality follows from Axiom~E4 (composition enriches):
$\rs(\mathrm{composeN}(a,n) \compose a, \mathrm{composeN}(a,n) \compose a) \geq
\rs(\mathrm{composeN}(a,n), \mathrm{composeN}(a,n))$.  Since $\mathrm{composeN}(a,n+1) =
\mathrm{composeN}(a,n) \compose a$, this gives
$\mathrm{weight}(\mathrm{composeN}(a,n+1)) \geq \mathrm{weight}(\mathrm{composeN}(a,n))$.
The second inequality follows from Axiom~R3 ($\rs(x,x) \geq 0$
for all $x$).
\end{proof}

\begin{lstlisting}
-- From IDT_v3b_DeepPhilosophy.lean

/-- DP22: Background Weight Grows Monotonically. -/
theorem background_weight_monotone (a : I) (n : ℕ) :
    weight (composeN a (n + 1)) ≥ weight (composeN a n) ∧
    weight (composeN a n) ≥ 0 := by
  unfold weight
  constructor
  · exact rs_composeN_mono a n
  · exact rs_composeN_nonneg a n
\end{lstlisting}

\begin{surprisebox}
\noindent\textbf{Why this is surprising.}
The monotonicity theorem says that the background \emph{never shrinks}.
Every idea you compose into the background stays there, permanently
increasing (or at least maintaining) the weight.  There is no mechanism
for forgetting, no mechanism for compression, no mechanism for reduction.
The frame problem is not a bug in our representational scheme; it is
a structural feature of ideatic space itself.  Even adding the void
idea---adding ``nothing''---preserves the weight (since $a \compose \void = a$),
and adding anything non-void \emph{increases} it.
\end{surprisebox}

\begin{interpretation}[Heidegger's Background Practices]
\label{interp:heidegger-background}
This theorem formalizes a central insight from Heidegger's analytic of
Dasein.  In \textit{Being and Time}, Heidegger argues that every act of
understanding occurs against a background of ``fore-having''
(\textit{Vorhabe}), ``fore-sight'' (\textit{Vorsicht}), and ``fore-conception''
(\textit{Vorgriff}).  These are not propositions stored in memory but
practical competences that structure the field of possibilities.

The monotonicity theorem captures the irreversibility of this background.
Once you have learned to ride a bicycle, you cannot \emph{un-learn} it.
Once you have understood a mathematical proof, the understanding becomes
part of your background---it colors every subsequent mathematical
encounter.  You can forget the details, but the structural change to
your ideatic configuration is permanent.

Dreyfus used this observation to argue that AI systems based on explicit
symbolic representations could never capture the background.  The
formal result here vindicates his intuition: the background is not
merely large but \emph{monotonically growing}.  No fixed-size
representation can capture it because it grows without bound.
\end{interpretation}

\section{Frame Explosion for Non-Void Ideas (DP23)}

\begin{theorem}[Frame Explosion --- DP23]
\label{thm:frame-explosion}
For any non-void idea $a \neq \void$ and any $n \in \N$:
\begin{enumerate}
\item $\mathrm{weight}(\mathrm{composeN}(a, n+1)) \geq \mathrm{weight}(a)$,
\item $\mathrm{weight}(\mathrm{composeN}(a, n+1)) > 0$,
\item $\mathrm{composeN}(a, n+1) \neq \void$.
\end{enumerate}
A single non-void idea generates an ever-present, strictly positive,
non-void frame at every stage of accumulation.
\end{theorem}

\begin{proof}
For (1): by induction on $n$.  The base case $n = 0$ gives
$\mathrm{composeN}(a, 1) = \void \compose a = a$, so the weight
equals $\mathrm{weight}(a)$.  The inductive step uses the monotonicity
from DP22.

For (3): we show $\mathrm{composeN}(a, n+1) \neq \void$ by proving that
composition with a non-void left factor cannot yield void.  Since
$\mathrm{composeN}(a, 1) = a \neq \void$, and each subsequent step
composes with $a$, the result is never void.

For (2): by Axiom~R4 (non-degeneracy), $x \neq \void$ implies
$\rs(x,x) > 0$.  Since $\mathrm{composeN}(a,n+1) \neq \void$ by (3),
we get $\mathrm{weight}(\mathrm{composeN}(a,n+1)) > 0$.
\end{proof}

\begin{lstlisting}
-- From IDT_v3b_DeepPhilosophy.lean

/-- DP23: Frame Explosion for Non-Void Ideas. -/
theorem frame_explosion (a : I) (h : a ≠ void) (n : ℕ) :
    weight (composeN a (n + 1)) ≥ weight a ∧
    weight (composeN a (n + 1)) > 0 ∧
    composeN a (n + 1) ≠ void := by
  have h1 : weight a > 0 := weight_pos a h
  have h2 : weight (composeN a (n + 1)) ≥ weight a :=
    rs_composeN_ge_base a n
  have h3 : composeN a (n + 1) ≠ void := composeN_ne_void a h n
  have h4 : weight (composeN a (n + 1)) > 0 := weight_pos _ h3
  exact ⟨h2, h4, h3⟩
\end{lstlisting}

\begin{surprisebox}
\noindent\textbf{Why this is surprising.}
A single idea, composed with itself $n$ times, generates a frame that is
\emph{at least as heavy as the idea itself} and \emph{strictly positive}.
This means that even the simplest possible non-void idea---a single
atomic concept---generates an unbounded frame.  The frame doesn't just
grow; it never collapses, never returns to void, and maintains at least
the weight of the original idea at every stage.

Compare this with the frame problem in AI: the difficulty is specifying
what doesn't change when an action is performed.  Here the mathematics
says something stronger: the set of things that \emph{do} change---the
frame itself---grows without bound.  It is not merely that the
unchanged things are hard to list; it is that the \emph{changed} things
proliferate endlessly.
\end{surprisebox}

\begin{interpretation}[The Chinese Room's Hidden Frame]
\label{interp:chinese-room}
John Searle's Chinese Room thought experiment (1980) imagines a person
in a room who manipulates Chinese symbols according to a rule book
without understanding Chinese.  Searle argues that the room does not
understand Chinese because symbol manipulation alone is insufficient
for understanding.

The frame explosion theorem reveals a hidden assumption in Searle's
argument.  The person in the room does have a frame---an enormous
background of English-language understanding, embodied skills, and
cultural knowledge---that is being used to follow the rules.  The
question is not whether the room ``understands Chinese'' but whether
the \emph{frame} of the rule-follower is adequate to the task.

In IDT terms, the rule-follower's frame grows with every symbol
manipulation.  After processing $n$ Chinese symbols, the frame has
weight at least $\mathrm{weight}(\text{symbol})$---the rules are
changing the rule-follower, even if Searle denies it.  Understanding
is not a binary state but a continuous accumulation of weight, and the
accumulation is irreversible.
\end{interpretation}

\section{The Incompressibility of Background (DP24)}

\begin{theorem}[Incompressible Background --- DP24]
\label{thm:incompressible-background}
For any ideas $a, c \in \IS$ and $n \in \N$:
\[
\emerge(\mathrm{composeN}(a,n),\, a,\, c)^2 \leq
\mathrm{weight}(\mathrm{composeN}(a, n+1)) \cdot \mathrm{weight}(c).
\]
The emergence from adding one more layer to the background is bounded by
the geometric mean of the accumulated background weight and the
observer's weight.
\end{theorem}

\begin{proof}
Since $\mathrm{composeN}(a,n+1) = \mathrm{composeN}(a,n) \compose a$,
this is a direct application of Axiom~E4 (emergence bounded):
$\emerge(x, y, c)^2 \leq \rs(x \compose y, x \compose y) \cdot \rs(c, c)$
with $x = \mathrm{composeN}(a,n)$ and $y = a$.
\end{proof}

\begin{lstlisting}
-- From IDT_v3b_DeepPhilosophy.lean

/-- DP24: Incompressible Background Emergence. -/
theorem incompressible_background (a c : I) (n : ℕ) :
    (emergence (composeN a n) a c) ^ 2 ≤
    weight (composeN a (n + 1)) * weight c := by
  unfold weight
  change (emergence (composeN a n) a c) ^ 2 ≤
    rs (compose (composeN a n) a) (compose (composeN a n) a) * rs c c
  exact emergence_bounded' (composeN a n) a c
\end{lstlisting}

\begin{surprisebox}
\noindent\textbf{Why this is surprising.}
The frame constrains its own growth rate.  While DP22 tells us the
background only grows, DP24 tells us it cannot grow \emph{arbitrarily
fast}.  The emergence from each new layer---the ``surprise'' of
adding one more instance of $a$ to the background---is bounded by the
square root of the background's current weight times the observer's
weight.  The richer the background becomes, the more room there is
for emergence, but the bound is sub-linear in the background weight.
The frame is like a snowball rolling downhill: it grows, but the
rate of growth relative to its size eventually decelerates.

This is the mathematical resolution of the frame problem: the
background is infinite and ever-growing, but it is \emph{not}
chaotic.  Its growth is lawful, bounded, and predictable.
\end{surprisebox}

\begin{interpretation}[Dreyfus and the Limits of AI]
\label{interp:dreyfus-limits}
Hubert Dreyfus spent decades arguing that artificial intelligence could
not replicate human intelligence because it could not capture the
background.  His argument, rooted in Heidegger and Merleau-Ponty, was
that the background is not a collection of facts but a holistic,
embodied, pre-reflective understanding that cannot be formalized.

The incompressibility theorem partially vindicates and partially
refines this argument.  Dreyfus was right that the background is
infinite and irreducible.  But he was wrong that it cannot be
formalized.  The ideatic axioms capture the background's behavior
precisely: it grows monotonically (DP22), it never collapses (DP23),
and its growth is bounded (DP24).  What cannot be done is to
\emph{list} the background explicitly.  But characterizing its
structural properties---which is what the axioms do---is entirely
possible.

The real lesson is not that the background is unformalizable but that
the \emph{kind} of formalization matters.  Explicit logical
representations (the approach Dreyfus criticized) cannot capture the
background because they require listing.  Algebraic characterizations
(the approach of IDT) can capture it because they describe structural
properties without enumeration.
\end{interpretation}

\section{The Frame Problem as Existential Condition}

Let us step back and consider the philosophical weight of these three
theorems together.

DP22 says the background only grows.  This is the mathematical form of
what Heidegger called ``thrownness'' (\textit{Geworfenheit}): we are
always already embedded in a background that we did not choose and
cannot escape.  Every encounter adds to the background, and nothing
subtracts from it.

DP23 says a single idea generates an unbounded frame.  This is the
mathematical form of what hermeneutic theorists call the
``fore-structure of understanding'': even the simplest act of
understanding presupposes an infinite context.  You cannot understand
the word ``dog'' without understanding language, biology, companionship,
loyalty, and an endless web of associations---each of which presupposes
further associations.

DP24 says the frame constrains its own growth.  This is the mathematical
form of what Gadamer called ``effective history''
(\textit{Wirkungsgeschichte}): the background is not an arbitrary
accumulation but a structured, lawful process whose growth is bounded by
its own internal logic.

Together, these three theorems constitute a complete mathematical
theory of the frame problem.  The frame is infinite (DP22--23) but
lawful (DP24).  It cannot be listed but can be characterized.  It
cannot be escaped but can be understood.

\section{The Frame Problem and the Limits of Logical AI}

The history of artificial intelligence can be read as a long struggle
with the frame problem.  The logical AI tradition---from McCarthy's
situation calculus to the modern knowledge representation
community---sought to solve the frame problem by explicitly listing
``frame axioms'': statements of the form ``action $A$ does not change
property $P$.''  The difficulty was combinatorial: for $n$ actions and
$m$ properties, you need $O(n \cdot m)$ frame axioms, and both $n$
and $m$ grow without bound.

Various technical solutions were proposed.  The ``closed world
assumption'' (Reiter, 1978) says that a property does not change
unless you explicitly assert that it does.  The ``Yale Shooting
Problem'' (Hanks and McDermott, 1987) showed that even this assumption
leads to paradoxes.  Nonmonotonic logics, circumscription, and
default reasoning were developed to handle the problem, each with
its own limitations.

The IDT perspective reveals why all these solutions are partial.
The frame is not a set of axioms but a \emph{weight}---a continuous,
monotonically increasing quantity that summarizes the accumulated
background.  No discrete set of frame axioms can capture a continuous
quantity.  The logical AI tradition was trying to discretize something
that is inherently continuous.

This does not mean that logical AI is useless.  It means that logical AI
operates in a regime where the frame problem is \emph{approximated}
by finite lists of axioms.  The approximation works in narrow domains
where the relevant background is small.  But in open-ended domains---natural
language understanding, common-sense reasoning, creative
problem-solving---the approximation breaks down because the frame is too
large and too dynamic.

\begin{theorem}[Frame Weight After $n$ Steps]
\label{thm:frame-weight-after-n}
For a non-void idea $a$ and any $n \geq 1$:
\[
\mathrm{weight}(\mathrm{composeN}(a, n)) \geq \mathrm{weight}(a) > 0.
\]
The frame weight at any step beyond the first is at least the weight of
the original idea, and strictly positive.
\end{theorem}

\begin{proof}
For $n = 1$: $\mathrm{composeN}(a, 1) = \void \compose a = a$, so
$\mathrm{weight}(\mathrm{composeN}(a, 1)) = \mathrm{weight}(a)$.
For $n > 1$: by DP22 (monotonicity), $\mathrm{weight}(\mathrm{composeN}(a, n))
\geq \mathrm{weight}(\mathrm{composeN}(a, 1)) = \mathrm{weight}(a)$.
Since $a \neq \void$, Axiom~R4 gives $\mathrm{weight}(a) > 0$.
\end{proof}

\begin{interpretation}[The Embodiment Thesis]
\label{interp:embodiment-thesis}
The frame problem is intimately connected to the \emph{embodiment thesis}
in philosophy of mind and AI.  Rodney Brooks, in his influential
``Intelligence Without Representation'' (1991), argued that intelligent
behavior does not require explicit symbolic representations of the world.
Instead, the world itself serves as its own model: the agent's body and
its interactions with the environment \emph{are} the background.

DP22--24 support a refined version of the embodiment thesis.  The
background is indeed too large for explicit representation (DP22--23).
But it is not unstructured---it obeys precise mathematical laws (DP24).
An embodied agent does not need to represent the frame explicitly
because the frame is carried in the agent's accumulated compositional
history.  The weight function is the mathematical summary of
everything the agent has encountered.

This is why deep learning systems, despite their lack of explicit
symbolic representations, can handle open-ended tasks that stymied
classical AI: they accumulate a frame through training that is
captured in their parameter weights, and these weights grow
monotonically (at least in expectation) during training.  The
parameters \emph{are} the frame.
\end{interpretation}

\section{The Cocycle Structure of the Frame}

The frame is not merely a weight; it has internal structure governed by
the cocycle condition.  The emergence from composing two ideas within
the frame satisfies:
\[
\emerge(a, b, d) + \emerge(a \compose b, c, d) =
\emerge(b, c, d) + \emerge(a, b \compose c, d).
\]
This constraint means that the frame's internal emergence patterns are
\emph{consistent}---they cannot be arbitrary but must satisfy a
global compatibility condition.

The cocycle condition connects the frame problem to cohomology theory.
In the language of group cohomology, the emergence function is a
$2$-cocycle on the monoid $(\IS, \compose, \void)$.  The frame is
therefore not just a weight but a \emph{cohomology class}---a global
invariant that captures the non-trivial interactions within the
accumulated background.

This observation has practical implications.  If the frame were
merely a weight, it could be summarized by a single number.  But
because the frame also carries cocycle structure, it contains richer
information about how its components interact.  The frame is not
a scalar but a \emph{function}---the emergence function restricted
to the accumulated background---and this function satisfies
algebraic constraints that give it global coherence.

\begin{reflectionbox}
\noindent\textbf{Reflection: The Frame Problem Is Not a Bug.}

The frame problem has traditionally been treated as an obstacle---something
to be overcome by clever engineering.  The theorems of this chapter suggest
a different perspective: the frame problem is a \emph{feature} of any
system rich enough to support genuine understanding.

If the background could be finitely listed, it would not be a real
background---it would be a finite set of propositions, and understanding
would reduce to logical deduction from a fixed database.  The frame
problem arises precisely because understanding is \emph{compositional}
and \emph{emergent}: each new composition generates new background, and
the background generates new possibilities for emergence.

The frame is not a bug.  It is the price of meaning.
\end{reflectionbox}

\section{Philosophical Coda}

The frame problem reveals a deep tension in the project of
artificial intelligence.  The early AI researchers---McCarthy, Minsky,
Newell, Simon---assumed that intelligence could be reduced to
symbol manipulation on a finite set of representations.  The frame
problem showed that this assumption was untenable: you cannot
represent all the consequences and non-consequences of an action
in a finite logical theory.

Dreyfus and the phenomenological tradition offered a diagnosis but
not a cure.  They showed that the background is holistic, embodied,
and pre-reflective, but they could not say anything precise about its
structure.  The result was a philosophical impasse: the AI researchers
knew the problem was real but had no formal tools to address it; the
phenomenologists had the right intuitions but no mathematics.

IDT breaks the impasse by providing a formal framework that
respects the phenomenological insights.  The background is infinite
(DP22--23), as Dreyfus insisted.  But it is also lawful (DP24), which
means it can be studied mathematically.  The key insight is that the
background is not a set of propositions but a \emph{weight}---a single
real number that summarizes the accumulated effect of all prior
compositions.  This weight grows monotonically, never collapses, and
constrains its own growth rate.

The frame problem is solved---not by listing the background
(which is impossible) but by characterizing its dynamics (which is
elegant).  The background is infinite, but its infinity is tame.

This, perhaps, is the deepest lesson of this chapter: the background
is not the enemy of understanding but its precondition.  Without the
ever-growing frame, there would be no context for emergence, no
structure for meaning, and no possibility of genuine understanding.
The frame problem is not a problem to be solved but a condition to
be embraced.


%%%%%%%%%%%%%%%%%%%%%%%%%%%%%%%%%%%%%%%%%%%%%%%%%%%%%%%%%%%%%%%%%%%%%%%%
%%  CHAPTER 7 — DIALECTICAL DIVERGENCE: WHEN HISTORY NEVER ENDS
%%%%%%%%%%%%%%%%%%%%%%%%%%%%%%%%%%%%%%%%%%%%%%%%%%%%%%%%%%%%%%%%%%%%%%%%

\chapter{Dialectical Divergence --- When History Never Ends}
\label{ch:dialectical-divergence}

\epigraph{``The history of all hitherto existing society is the history of
class struggles.  Freeman and slave, patrician and plebeian, lord and serf
\ldots\ in a word, oppressor and oppressed, stood in constant opposition
to one another.''}{Karl Marx, \textit{The Communist Manifesto} (1848)}

\section{The Hegelian Dialectic and Its Fate}

Georg Wilhelm Friedrich Hegel proposed the most ambitious philosophy of
history ever conceived.  History, he argued, is the progressive
self-realization of Spirit (\textit{Geist}) through a dialectical process:
thesis confronts antithesis, producing synthesis, which becomes the new
thesis for the next round.  Each stage incorporates (\textit{aufhebt})
its predecessor---simultaneously preserving, negating, and elevating it.
The process converges, Hegel claimed, toward Absolute Knowledge: a state
in which Spirit fully comprehends itself.

Marx materialized the dialectic, replacing Spirit with class struggle.
History is driven by contradictions in the mode of production: feudalism
generates the bourgeoisie, capitalism generates the proletariat, and the
proletariat's revolution produces a classless society.  Marx shared
Hegel's conviction that the dialectical process has an end point---the
communist society in which all contradictions are resolved.

In 1992, Francis Fukuyama provocatively declared ``the end of history'':
with the fall of the Soviet Union and the global triumph of liberal
democracy, the dialectical process had reached its terminus.  History
had achieved its final synthesis.

In this chapter, we prove that all three thinkers are mathematically
wrong---at least within the ideatic space.  The dialectical process
does not converge.  History never ends.  The only way for the dialectic
to terminate is for it to have started from nothing.

\section{Iterated Dialectical Synthesis}

\begin{definition}[Iterated Dialectic]
\label{def:iter-dialectic}
Given a \textbf{thesis} $a \in \IS$ and an \textbf{antithesis} $b \in \IS$,
the \textbf{iterated dialectical synthesis} is defined recursively:
\begin{align}
\mathrm{iterDialectic}(a, b, 0) &= a, \\
\mathrm{iterDialectic}(a, b, n+1) &=
\mathrm{iterDialectic}(a, b, n) \compose b.
\end{align}
\end{definition}

At stage $0$, we have the thesis alone.  At stage $1$, the thesis
encounters the antithesis: $a \compose b$.  At stage $2$, the
synthesis encounters the antithesis again: $(a \compose b) \compose b$.
And so on.  Each stage produces a new synthesis that incorporates all
previous encounters.

Note that this definition is asymmetric: the antithesis is always composed
on the right.  This captures the Hegelian intuition that the antithesis is
the \emph{challenge}---the external opposition that forces the thesis to
develop.  The thesis grows by absorbing the antithesis, stage after stage,
in an endless dialectical spiral.

\begin{lstlisting}
-- From IDT_v3b_DeepPhilosophy.lean

/-- Iterated dialectical synthesis. -/
def iterDialectic (thesis antithesis : I) : ℕ → I
  | 0 => thesis
  | n + 1 => compose (iterDialectic thesis antithesis n) antithesis
\end{lstlisting}

\section{History Accumulates (DP25)}

\begin{theorem}[History Accumulates Weight --- DP25]
\label{thm:history-accumulates}
For any thesis $a$, antithesis $b$, and stage $n$:
\[
\mathrm{weight}(\mathrm{iterDialectic}(a, b, n+1)) \geq
\mathrm{weight}(\mathrm{iterDialectic}(a, b, n)) \geq 0.
\]
The weight of the dialectical synthesis is non-decreasing and always
non-negative.  History never retreats.
\end{theorem}

\begin{proof}
By Axiom~E4 (composition enriches):
$\rs(\mathrm{iterDialectic}(a,b,n) \compose b,\,
\mathrm{iterDialectic}(a,b,n) \compose b) \geq
\rs(\mathrm{iterDialectic}(a,b,n),\, \mathrm{iterDialectic}(a,b,n))$.
This gives the first inequality.  The second follows from
Axiom~R3 ($\rs(x,x) \geq 0$).
\end{proof}

\begin{lstlisting}
-- From IDT_v3b_DeepPhilosophy.lean

/-- DP25: History Accumulates Weight Monotonically. -/
theorem history_accumulates (a b : I) (n : ℕ) :
    weight (iterDialectic a b (n + 1)) ≥ weight (iterDialectic a b n) ∧
    weight (iterDialectic a b n) ≥ 0 := by
  unfold weight
  constructor
  · exact compose_enriches' (iterDialectic a b n) b
  · exact rs_self_nonneg' _
\end{lstlisting}

\begin{surprisebox}
\noindent\textbf{Why this is surprising.}
Even when thesis and antithesis are bitterly opposed---when their mutual
resonance is deeply negative---the synthesis is \emph{always heavier}
than the thesis alone.  Opposition does not diminish; it enriches.
This runs counter to the naive intuition that conflict is destructive.
In the ideatic space, conflict is always productive: the clash of
ideas generates weight, never loses it.

Consider what this means for Hegel: the dialectical process cannot
``go backward.''  You cannot have a stage of history that is lighter
than a previous stage.  The synthesis always incorporates the thesis,
and incorporation always adds weight.  History is a ratchet: it clicks
forward but never clicks back.
\end{surprisebox}

\begin{interpretation}[The Irreversibility of Cultural Change]
\label{interp:cultural-irreversibility}
This theorem captures something that historians and anthropologists
have long observed: cultural change is irreversible.  You cannot
``un-invent'' the printing press, ``un-discover'' evolution, or
``un-think'' the concept of human rights.  Once an idea enters the
cultural synthesis, it becomes part of the background, and the
background only grows (DP22).

The Romantics tried to return to a pre-industrial pastoral; they
failed.  The Luddites tried to destroy the machines; the machines
persisted.  Every reactionary movement in history has discovered
that you cannot turn back the clock: the best you can do is compose
the past with the present, producing a synthesis that is heavier
than either.

Marx understood this intuitively when he distinguished between
\textit{Aufhebung}---the dialectical ``lifting up'' that simultaneously
preserves and negates---and simple destruction.  The formal result
shows that \textit{Aufhebung} is the \emph{only} possibility: the
synthesis always preserves the weight of the thesis, because
composition always enriches.
\end{interpretation}

\section{The End of History Requires Void (DP26)}

\begin{theorem}[End of History Requires Void --- DP26]
\label{thm:end-of-history}
If the thesis is non-void ($a \neq \void$), then for all $n \in \N$:
\[
\mathrm{iterDialectic}(a, b, n) \neq \void.
\]
The dialectical process can only reach void (silence, the absence of all
ideas) if it started from void.  History can only ``end'' if it never began.
\end{theorem}

\begin{proof}
By induction on $n$.  Base case: $\mathrm{iterDialectic}(a,b,0) = a \neq \void$.
Inductive step: suppose $\mathrm{iterDialectic}(a,b,n) \neq \void$.  Then
$\mathrm{iterDialectic}(a,b,n+1) = \mathrm{iterDialectic}(a,b,n) \compose b$.
Since the left factor $\mathrm{iterDialectic}(a,b,n)$ is non-void,
the composition is non-void (by the derived lemma
\texttt{compose\_ne\_void\_of\_left}).
\end{proof}

\begin{lstlisting}
-- From IDT_v3b_DeepPhilosophy.lean

/-- DP26: The End of History Requires Void. -/
theorem end_of_history_requires_void (a b : I) (h : a ≠ void) (n : ℕ) :
    iterDialectic a b n ≠ void := by
  induction n with
  | zero => exact h
  | succ n ih =>
    show compose (iterDialectic a b n) b ≠ void
    exact compose_ne_void_of_left _ b ih
\end{lstlisting}

\begin{surprisebox}
\noindent\textbf{Why this is surprising.}
Fukuyama's ``end of history'' is mathematically impossible for any
non-void thesis.  This is not a historical or sociological argument;
it is a \emph{theorem}.  As long as the initial thesis has positive
weight---as long as there is \emph{something} rather than nothing---the
dialectical process never reaches void.  History never falls silent.

Even more striking: the theorem holds regardless of what the antithesis
is.  You can oppose the thesis with its perfect negation, with chaos,
with anything at all---the synthesis is never void.  The dialectical
process is \emph{robustly} non-terminating.

Hegel imagined history converging to Absolute Knowledge.  Marx imagined
it converging to the classless society.  Fukuyama imagined it
converging to liberal democracy.  The theorem says: it converges to
nothing, because it never stops.
\end{surprisebox}

\begin{interpretation}[Fukuyama Refuted]
\label{interp:fukuyama-refuted}
In 1992, Francis Fukuyama published \textit{The End of History and the
Last Man}, arguing that the worldwide spread of liberal democracy
marked the ``end of history''---the final stage of Hegel's dialectical
process.  The argument was immediately controversial, and subsequent
events (9/11, the rise of China, the resurgence of authoritarianism)
seemed to refute it empirically.

DP26 provides a mathematical refutation.  For history to ``end'' in
the IDT sense---for the dialectical process to reach void---the thesis
must have been void to begin with.  Since liberal democracy is
manifestly non-void (it has positive weight as a political idea),
the dialectical process that produced it cannot terminate.  New
antitheses will always arise, new syntheses will always form, and
the weight of the cultural accumulation will always grow.

This does not mean that history is chaotic or directionless.  DP25
says that the weight is monotonically increasing, which implies a
kind of progress---but it is progress in \emph{weight}, not
necessarily in justice, freedom, or human flourishing.  The dialectic
accumulates; whether it improves is a separate question that the
axioms do not address.
\end{interpretation}

\section{Dialectical Weight Is Globally Monotone (DP27)}

\begin{theorem}[Dialectical Weight Is Transitive --- DP27]
\label{thm:dialectical-transitive}
For any thesis $a$, antithesis $b$, and stages $m \leq n$:
\[
\mathrm{weight}(\mathrm{iterDialectic}(a, b, n)) \geq
\mathrm{weight}(\mathrm{iterDialectic}(a, b, m)).
\]
The weight ordering of dialectical stages is globally monotone, not just
locally.
\end{theorem}

\begin{proof}
By induction on $n$.  If $m = n$, the result is immediate.  If $m < n+1$,
then $m \leq n$, so by the inductive hypothesis,
$\mathrm{weight}(\mathrm{iterDialectic}(a,b,n)) \geq
\mathrm{weight}(\mathrm{iterDialectic}(a,b,m))$.  By DP25,
$\mathrm{weight}(\mathrm{iterDialectic}(a,b,n+1)) \geq
\mathrm{weight}(\mathrm{iterDialectic}(a,b,n))$.  Transitivity of $\geq$
gives the result.
\end{proof}

\begin{lstlisting}
-- From IDT_v3b_DeepPhilosophy.lean

/-- DP27: Dialectical Weight Is Transitive. -/
theorem dialectical_weight_transitive (a b : I) (m n : ℕ) (h : m ≤ n) :
    weight (iterDialectic a b n) ≥ weight (iterDialectic a b m) := by
  unfold weight
  induction n with
  | zero =>
    interval_cases m
    exact le_refl _
  | succ n ih =>
    rcases Nat.eq_or_lt_of_le h with rfl | hlt
    · exact le_refl _
    · have h1 := ih (Nat.lt_succ_iff.mp hlt)
      have h2 : rs (iterDialectic a b (n + 1))
                    (iterDialectic a b (n + 1)) ≥
                rs (iterDialectic a b n)
                    (iterDialectic a b n) :=
        compose_enriches' _ b
      linarith
\end{lstlisting}

\begin{surprisebox}
\noindent\textbf{Why this is surprising.}
DP25 gives local monotonicity: each stage is at least as heavy as its
immediate predecessor.  DP27 extends this to \emph{global} monotonicity:
any later stage is at least as heavy as any earlier stage, no matter how
many stages separate them.  This seems like it should be obvious---if
each step goes up, then the whole sequence goes up---but the formal
proof requires a careful induction that handles the base case and the
transition from local to global.

The philosophical import is that you cannot ``skip back'' in history.
Stage $100$ of the dialectic is guaranteed to be at least as heavy as
stage $3$.  You can leap forward but never backward.  The arrow of
history points in one direction only: toward greater weight.
\end{surprisebox}

\begin{interpretation}[Marx's Historical Materialism Formalized]
\label{interp:marx-formalized}
Marx's historical materialism holds that history progresses through
stages---from primitive communism to slavery to feudalism to capitalism
to socialism.  Each stage is ``higher'' than the previous one in the
sense that it involves more developed productive forces and more
complex social relations.

DP27 formalizes the progressive aspect of this view.  The weight
function provides a precise measure of ``development'': later
dialectical stages are guaranteed to have at least as much weight as
earlier ones.  Marx was right that history doesn't go backward---but
the formal result is even stronger.  Marx thought only material
conditions were irreversible; the theorem shows that the \emph{ideatic}
development is equally irreversible.  Ideas, once synthesized, cannot
be un-synthesized.

However, the theorem also reveals a crucial limitation of Marx's
view.  The monotonicity is in \emph{weight}, not in any normative
measure.  History may accumulate weight without accumulating justice.
The feudal synthesis may be heavier than the slave synthesis without
being morally superior.  The theorem captures the inevitability of
change but is silent on its direction.
\end{interpretation}

\section{The Dialectical Spiral}

Together, DP25--27 paint a picture of the dialectical process that
diverges sharply from both Hegel and Marx.

For Hegel, the dialectic \emph{converges}: it spirals inward toward
Absolute Knowledge, where Spirit fully comprehends itself.  The
formal result says the opposite.  The dialectic \emph{diverges}: it
spirals outward, accumulating weight without bound, never reaching a
fixed point (unless it started from void).

For Marx, the dialectic has a \emph{terminus}: the classless society
in which all contradictions are resolved.  The formal result says
there is no terminus.  New contradictions are always possible because
each synthesis has strictly positive weight and therefore generates
new emergence when composed with anything.

The picture that emerges is closer to what Walter Benjamin described
in his ``Theses on the Philosophy of History'' (1940): the angel of
history, blown backward by the storm of progress, watching the
wreckage pile higher and higher.  The dialectical process is
irreversible, cumulative, and open-ended.  It does not converge to
utopia; it diverges into complexity.

\section{The Emergence of the Dialectical Process}

The dialectical process generates emergence at every stage.  When the
synthesis at stage $n$ encounters the antithesis $b$ to produce stage
$n+1$, the emergence is:
\[
\emerge(\mathrm{iterDialectic}(a,b,n),\, b,\, d)
= \rs(\mathrm{iterDialectic}(a,b,n+1),\, d)
- \rs(\mathrm{iterDialectic}(a,b,n),\, d) - \rs(b,\, d).
\]
This emergence captures the genuinely new content generated by each
dialectical step---the part of the synthesis that was not present
in either the thesis or the antithesis alone.

By the Cauchy--Schwarz bound (Axiom~E4), this emergence is bounded:
\[
\emerge(\mathrm{iterDialectic}(a,b,n),\, b,\, d)^2 \leq
\mathrm{weight}(\mathrm{iterDialectic}(a,b,n+1)) \cdot \mathrm{weight}(d).
\]
The dialectical process produces emergence at every stage, but the
emergence at each stage is bounded by the accumulated weight of
the synthesis.  The richer the history, the more room for emergence
at the next step---but the room grows only as the square root.

\begin{interpretation}[Benjamin's Angel of History]
\label{interp:benjamin-angel}
Walter Benjamin's ninth thesis on the philosophy of history describes
Paul Klee's painting \textit{Angelus Novus}:

\begin{quote}
His face is turned toward the past.  Where we perceive a chain of
events, he sees one single catastrophe which keeps piling wreckage upon
wreckage and hurls it in front of his feet.  The angel would like to
stay, awaken the dead, and make whole what has been smashed.  But a
storm is blowing from Paradise; it has got caught in his wings with such
violence that the angel can no longer close them.  This storm
irresistibly propels him into the future to which his back is turned,
while the pile of debris before him grows skyward.
\end{quote}

The mathematics of DP25--27 gives Benjamin's image a precise form.  The
``pile of debris'' is the monotonically increasing weight of the
dialectical synthesis.  The ``storm from Paradise'' is the axiom of
composition enrichment, which ensures that each new composition adds
weight.  The angel cannot ``make whole what has been smashed'' because
the weight function is monotone: there is no mechanism for returning
to a lighter state.

Benjamin saw the dialectic as catastrophe; Hegel saw it as progress.
The mathematics is neutral: it says only that the weight grows.
Whether the growth is catastrophic or progressive depends on what
the weight measures---and the axioms do not tell us that.
\end{interpretation}

\section{The Dialectical Process and Thermodynamic Irreversibility}

There is a deep analogy between the dialectical weight function and
entropy in thermodynamics.  The second law of thermodynamics states
that the entropy of an isolated system never decreases.  DP25 states
that the weight of a dialectical synthesis never decreases.  Both
are statements of irreversibility: the system can only go ``forward''
in a well-defined sense.

The analogy extends further.  In thermodynamics, entropy measures the
number of microstates consistent with the macrostate.  In IDT, weight
measures the ``ideatic content'' of the synthesis---the richness of
its compositional structure.  Just as a high-entropy state is one
with many microstates, a high-weight synthesis is one with a rich
internal structure.

The second law of thermodynamics is a consequence of the overwhelmingly
large number of microstates that correspond to high-entropy macrostates.
The monotonicity of dialectical weight is a consequence of the axiom
of composition enrichment, which ensures that composition never
reduces self-resonance.

Both laws have the same philosophical implication: the arrow of time
is real.  You cannot un-stir the cream from the coffee, and you cannot
un-synthesize the dialectical product.  The past is irrecoverable,
not merely in practice but in principle.

\begin{reflectionbox}
\noindent\textbf{Reflection: Progress Without Telos.}

The theorems of this chapter prove that history is progressive in a
strictly formal sense: later stages are always at least as heavy as
earlier stages.  But this is progress \emph{without telos}---accumulation
without a goal.  The dialectic grows heavier, but nothing in the axioms
says it grows \emph{better}.

This is, perhaps, the most honest account of history that mathematics
can provide.  History is irreversible, cumulative, and productive.
Whether it is also \emph{good} is a question that lies beyond the reach
of any formal theory.

The ideatic space gives us the arrow of time but not its meaning.
That meaning---if it exists---must come from outside the axioms.
\end{reflectionbox}

\section{Philosophical Coda}

The dialectical divergence theorems challenge every major philosophy
of history that posits a terminal state.  Hegel's Absolute Knowledge,
Marx's communist utopia, Fukuyama's liberal democratic paradise---all
require the dialectic to terminate, and DP26 proves that termination
requires void.

But the theorems are not nihilistic.  They do not say that history is
meaningless; they say that history is \emph{inexhaustible}.  Every
synthesis generates new possibilities for emergence, new tensions with
new antitheses, new stages of development.  The dialectic is not a
circle that closes but a spiral that opens.

This may be the most important philosophical conclusion of the entire
volume: the process of ideatic development is inherently open-ended.
Ideas compose, generate emergence, accumulate weight, and produce new
ideas---and this process has no natural stopping point.

The ``end of history'' is not just empirically unlikely; it is
mathematically impossible.  As long as there are ideas---as long as
the thesis is non-void---the dialectic continues.  The only silence
is the silence that was there before the first idea.  Everything after
that is noise, accumulation, and the irreversible growth of meaning.


%%%%%%%%%%%%%%%%%%%%%%%%%%%%%%%%%%%%%%%%%%%%%%%%%%%%%%%%%%%%%%%%%%%%%%%%
%%  CHAPTER 8 — THE ZOMBIE THEOREM: CONSCIOUSNESS LEAVES NO ROOM
%%                FOR ZOMBIES
%%%%%%%%%%%%%%%%%%%%%%%%%%%%%%%%%%%%%%%%%%%%%%%%%%%%%%%%%%%%%%%%%%%%%%%%

\chapter{The Zombie Theorem --- Consciousness Leaves No Room for Zombies}
\label{ch:zombie-theorem}

\epigraph{``Can we conceive of a being physically identical to a conscious
being, molecule for molecule, but lacking consciousness entirely?  If so,
then consciousness is something over and above the physical.
\ldots\ The conceivability of zombies seems to establish the conclusion that
consciousness is not logically supervenient on the physical.''}{David Chalmers, \textit{The Conscious Mind} (1996)}

\section{The Hard Problem of Consciousness}

The \emph{hard problem of consciousness}, as formulated by David Chalmers
in 1995, is this: why is there something it is like to be conscious?
Why does the neural processing of visual information feel like anything
at all?  Why isn't the brain just a complex information-processing
machine that operates ``in the dark,'' with no inner light of experience?

The ``easy problems'' of consciousness---explaining discrimination,
integration, reportability, attention, behavioral control---are
problems of functional organization.  They are hard in practice but
easy in principle: we can see, at least in outline, how a physical
system could perform these functions.  The hard problem is different
in kind.  Even if we explained every functional capacity of the brain,
we would still face the question: why is there \emph{experience}?

Chalmers's most celebrated argument for the irreducibility of
consciousness is the \emph{zombie argument}.  A philosophical zombie
(p-zombie) is a being that is physically identical to a conscious being
but has no conscious experience whatsoever.  It behaves exactly like you,
reports experiences just like you, and processes information just like
you---but there is ``nothing it is like'' to be a zombie.  If zombies
are \emph{conceivable}---if there is no logical contradiction in their
description---then consciousness cannot be logically entailed by
physical structure alone.  Consciousness must be something \emph{extra}.

In this chapter, we prove that within the ideatic space, zombies are
\emph{impossible}.  Not merely inconceivable but mathematically
ruled out.  If two ideas have the same emergence profile---the same
functional role in every compositional context---then they are
identical in every way that matters.  There is nothing left for a
zombie to lack.

\section{Same Emergence as Functional Equivalence}

\begin{definition}[Same Idea]
\label{def:same-idea}
Two ideas $a, b \in \IS$ are \textbf{functionally equivalent}
(written $\mathrm{sameIdea}(a, b)$) if and only if they have the
same emergence in every context:
\[
\mathrm{sameIdea}(a, b) \iff \forall c, d \in \IS:\,
\emerge(a, c, d) = \emerge(b, c, d).
\]
\end{definition}

This is a strong condition.  It says that $a$ and $b$ produce the
same emergence when composed with \emph{any} idea and observed by
\emph{any} idea.  They are indistinguishable from the standpoint of
compositional behavior.  If consciousness is supposed to be
``something extra'' beyond functional role, then sameIdea pairs are
precisely the candidates for zombie-hood: they do everything the
same but might (allegedly) differ in experience.

\begin{lstlisting}
-- From IDT_v3b_DeepPhilosophy.lean

/-- Two ideas have the same emergence profile. -/
def sameIdea (a b : I) : Prop :=
  ∀ c d : I, emergence a c d = emergence b c d
\end{lstlisting}

\section{Zombie Impossibility (DP28)}

\begin{theorem}[Zombie Impossibility --- DP28]
\label{thm:zombie-impossibility}
If $\mathrm{sameIdea}(a, b)$, then for all $c, d \in \IS$:
\[
\rs(a \compose c,\, d) - \rs(a,\, d) = \rs(b \compose c,\, d) - \rs(b,\, d).
\]
Ideas with the same emergence profile produce the same compositional
shift in resonance.  Their ``functional contribution'' to any composition
is identical.
\end{theorem}

\begin{proof}
By definition of emergence:
$\emerge(a, c, d) = \rs(a \compose c, d) - \rs(a, d) - \rs(c, d)$.
If $\emerge(a, c, d) = \emerge(b, c, d)$, then:
\[
\rs(a \compose c, d) - \rs(a, d) - \rs(c, d) =
\rs(b \compose c, d) - \rs(b, d) - \rs(c, d).
\]
Canceling $\rs(c, d)$ from both sides gives the result.
\end{proof}

\begin{lstlisting}
-- From IDT_v3b_DeepPhilosophy.lean

/-- DP28: Functional Equivalence Is Complete. -/
theorem zombie_impossibility (a b : I) (h : sameIdea a b) (c d : I) :
    rs (compose a c) d - rs a d = rs (compose b c) d - rs b d := by
  have h1 := h c d
  unfold emergence at h1
  linarith
\end{lstlisting}

\begin{surprisebox}
\noindent\textbf{Why this is surprising.}
This theorem dissolves the hard problem of consciousness within the
ideatic framework.  Chalmers argues that functional organization
(emergence profile) might not suffice for consciousness---that there
might be ``something extra'' beyond functional role.  DP28 proves that
in the ideatic space, there is \emph{nothing extra}.

If two ideas have the same emergence profile, then they shift
resonance by exactly the same amount in every composition.  Their
``compositional contribution''---the change they induce in any
context---is identical.  There is no room for a zombie because
there is no room for a difference that doesn't show up in some
compositional context.

The zombie is not merely inconceivable; it is \emph{algebraically
impossible}.  The axioms of the ideatic space leave no gap between
function and experience.
\end{surprisebox}

\begin{interpretation}[Dissolving the Hard Problem]
\label{interp:dissolving-hard-problem}
The zombie argument depends on a crucial assumption: that there is a
gap between functional organization and conscious experience.  Chalmers
calls this the ``explanatory gap''---the gap between third-person
descriptions of brain function and first-person descriptions of
experience.

DP28 closes this gap within the ideatic space.  The emergence
profile---which captures the functional role of an idea in every
compositional context---fully determines its compositional behavior.
If two ideas have the same emergence profile, they are functionally
identical in every sense that the axioms can detect.

This does not mean that consciousness is ``nothing but'' function.
It means that in the ideatic space, consciousness \emph{is} the
emergence profile.  Experience is not something added on top of
function; it is the function itself, viewed from the inside.
The ``hard problem'' dissolves because the distinction between
function and experience was never well-defined in the first place.

This is close to what Daniel Dennett has argued for decades: that
once you have explained all the functional capacities of the brain,
you have explained everything.  But where Dennett's argument is
philosophical and persuasive, DP28 is mathematical and demonstrative.
The axioms \emph{force} the conclusion.
\end{interpretation}

\section{No Hidden Qualia (DP29)}

\begin{theorem}[No Hidden Qualia --- DP29]
\label{thm:no-hidden-qualia}
If $\mathrm{sameIdea}(a, b)$, then for all $c, d \in \IS$:
\[
\rs(a \compose c,\, d) - \rs(b \compose c,\, d) = \rs(a,\, d) - \rs(b,\, d).
\]
The only difference between two functionally equivalent ideas is their
individual resonance---their ``denotation'' or ``connotation.''  There
are no hidden qualia.
\end{theorem}

\begin{proof}
From the proof of DP28:
$\rs(a \compose c, d) - \rs(a, d) = \rs(b \compose c, d) - \rs(b, d)$.
Rearranging:
$\rs(a \compose c, d) - \rs(b \compose c, d) = \rs(a, d) - \rs(b, d)$.
\end{proof}

\begin{lstlisting}
-- From IDT_v3b_DeepPhilosophy.lean

/-- DP29: No Hidden Qualia. -/
theorem no_hidden_qualia (a b : I) (h : sameIdea a b) (c d : I) :
    rs (compose a c) d - rs (compose b c) d =
    rs a d - rs b d := by
  have h1 := h c d
  unfold emergence at h1
  linarith
\end{lstlisting}

\begin{surprisebox}
\noindent\textbf{Why this is surprising.}
If two ideas have the same emergence profile, then the only way they
can differ is in their individual resonance values $\rs(a,d)$ vs.\
$\rs(b,d)$.  This individual resonance is the ``raw denotation'' of
the idea---how it resonates with other ideas in isolation, outside
of any compositional context.

This means there are exactly two components to an idea's identity:
\begin{enumerate}
\item Its \emph{emergence profile}: how it contributes to compositions
      (captured by $\emerge(a, c, d)$ for all $c, d$).
\item Its \emph{resonance profile}: how it relates to other ideas
      individually (captured by $\rs(a, d)$ for all $d$).
\end{enumerate}
There is no third component.  No hidden qualia, no secret inner
light, no ``what it is like'' that escapes these two descriptions.
The idea is fully characterized by its emergence and resonance.
\end{surprisebox}

\begin{interpretation}[Jackson's Mary's Room]
\label{interp:marys-room}
Frank Jackson's ``Mary's Room'' thought experiment (1982) imagines a
brilliant neuroscientist, Mary, who has spent her entire life in a
black-and-white room.  She knows everything there is to know about the
physics and neuroscience of color perception.  Then she steps outside
and sees red for the first time.  Does she learn something new?

Jackson argued yes: Mary gains knowledge of \emph{what it is like} to
see red---knowledge that cannot be derived from physical facts alone.
This is supposed to show that there are non-physical facts about
consciousness.

DP29 offers a different analysis.  Before leaving the room, Mary has
a certain resonance profile: $\rs(\mathrm{Mary}_{\mathrm{before}}, d)$
for all $d$.  After seeing red, she has a new resonance profile:
$\rs(\mathrm{Mary}_{\mathrm{after}}, d)$.  The ``new knowledge'' Mary
gains is not a hidden quale but a change in her resonance
profile---specifically, the emergence $\emerge(\mathrm{Mary}_{\mathrm{before}},
\mathrm{red}, d)$ that arises from composing her prior knowledge with
the experience of seeing red.

Mary does learn something new, but what she learns is \emph{emergence},
not qualia.  The emergence is fully determined by the axioms; there is
nothing mysterious about it.  The ``aha moment'' of seeing red is the
emergence function evaluated at specific arguments---a real number, not
a metaphysical mystery.
\end{interpretation}

\section{Same Emergence, Same Experience (DP30)}

\begin{theorem}[Same Emergence, Same Experience --- DP30]
\label{thm:same-emergence-experience}
If $\mathrm{sameIdea}(a, b)$, then for all $c, d \in \IS$:
\[
\emerge(a, c, d) = \emerge(b, c, d).
\]
Ideas with the same emergence profile have the same experience of
every encounter.
\end{theorem}

\begin{proof}
This is immediate from the definition of $\mathrm{sameIdea}(a, b)$,
which asserts exactly this condition for all $c$ and $d$.
\end{proof}

\begin{lstlisting}
-- From IDT_v3b_DeepPhilosophy.lean

/-- DP30: Same Emergence Implies Same Self-Composition Surplus. -/
theorem same_emergence_same_experience (a b : I) (h : sameIdea a b)
    (c d : I) : emergence a c d = emergence b c d := h c d
\end{lstlisting}

\begin{surprisebox}
\noindent\textbf{Why this is surprising.}
DP30 is almost tautological---it follows directly from the definition
of sameIdea.  But its philosophical significance is enormous.  It says
that the \emph{experience} of composition---the emergence that arises
when an idea meets another idea---is fully captured by the emergence
profile.  If two ideas have the same emergence profile, they have the
same experience in every encounter.

There is no ``inner light'' of consciousness that exists beyond the
emergence profile.  The experience IS the emergence.  This is the
strongest possible form of functionalism: not just that function
determines behavior, but that function \emph{is} experience.

The Lean proof is a single function application: \texttt{h c d}.
The entire argument for the identity of function and experience
reduces to applying the hypothesis.  The axioms leave no room for
anything else.
\end{surprisebox}

\begin{interpretation}[Nagel's Bat]
\label{interp:nagels-bat}
Thomas Nagel's famous essay ``What Is It Like to Be a Bat?'' (1974)
argues that there is a subjective character of experience---a ``what
it is like''---that cannot be captured by any objective, third-person
description.  No matter how much we know about bat sonar, we can never
know what it is \emph{like} to echolocate.

DP28--30 offer a precise response.  In the ideatic space, ``what it is
like'' to be an idea $a$ encountering an idea $c$ in the presence of
$d$ is \emph{exactly} $\emerge(a, c, d)$.  This is a real number---a
perfectly objective, third-person quantity.  If another idea $b$ has
the same emergence profile as $a$, then it has the same ``what it is
like'' in every context.  There is nothing left over.

Nagel would object that the emergence function cannot capture the
subjective \emph{feel} of experience.  But the theorems show that
within the ideatic space, there is no ``feel'' beyond the emergence.
The emergence determines all compositional behavior (DP28), determines
all differential behavior (DP29), and constitutes the entirety of
the experiential profile (DP30).

The ``what it is like'' is not a mysterious extra ingredient.  It is
the emergence function, evaluated at specific arguments.  Consciousness
is not a substance; it is a function.
\end{interpretation}

\section{The Identity Theory of Mind and Emergence}

The theorems of this chapter support a strong version of the
\emph{identity theory} of mind, but with a crucial twist.  Classical
identity theory (Smart, Place) identifies mental states with brain
states.  Functionalism (Putnam, early) identifies mental states with
functional roles.  The IDT version identifies mental states with
\emph{emergence profiles}.

The emergence profile of an idea $a$ is the function
$\emerge(a, \cdot, \cdot) : \IS \times \IS \to \R$.  This function
completely determines how $a$ contributes to compositions (DP28) and
how it differs from any other idea (DP29).  Two ideas with the same
emergence profile are, for all compositional purposes, the same idea
(DP30).

This is stronger than classical functionalism because emergence is
not merely ``causal role'' but \emph{algebraic structure}.  The
emergence function is constrained by the cocycle condition
(Axiom~E4), by the Cauchy--Schwarz bound (also Axiom~E4), and by
the interplay between composition and resonance.  It is not any
arbitrary function but a highly structured one.

\section{The Spectrum of Consciousness}

An immediate consequence of the zombie theorems is that consciousness
in the ideatic space is not binary but \emph{graded}.  Two ideas can
have emergence profiles that are ``close'' without being identical.
This suggests a natural metric on consciousness:

\begin{definition}[Emergence Distance]
\label{def:emergence-distance}
The \textbf{emergence distance} between two ideas $a, b \in \IS$ is:
\[
d_{\emerge}(a, b) := \sup_{c, d \in \IS}
|\emerge(a, c, d) - \emerge(b, c, d)|.
\]
Two ideas are functionally equivalent ($\mathrm{sameIdea}(a,b)$) if
and only if $d_{\emerge}(a, b) = 0$.
\end{definition}

This distance measures how ``different'' two ideas are in their
experiential profiles.  Ideas with small emergence distance are
``almost the same'' in their compositional behavior---they are
near-zombies of each other, differing only slightly in their
experience of the world.

The zombie theorems (DP28--30) say that $d_{\emerge}(a,b) = 0$
implies complete functional equivalence.  But the converse is
also interesting: if $d_{\emerge}(a,b) > 0$, then $a$ and $b$
\emph{must} differ in some compositional context.  There is no
way to have a non-zero emergence distance without a detectable
functional difference.

This dissolves the ``spectrum problem'' in consciousness studies:
the question of where consciousness begins and ends.  In the IDT
framework, there is no sharp boundary.  Every idea has an emergence
profile, and the richness of that profile determines the richness
of its ``experience.''  The void has a trivial emergence profile
($\emerge(\void, c, d) = 0$ for all $c, d$), and non-void ideas
have non-trivial emergence profiles.  Consciousness is not a
binary switch but a continuous variation in the structure of
emergence.

\section{Objections and Replies}

\subsection{The Knowledge Argument Revisited}

One might object that DP28--30 prove too much.  If consciousness is
nothing but emergence, then Mary in her black-and-white room should
be able to predict the emergence she will experience upon seeing red,
since she knows ``everything there is to know'' about color perception.
Does she not already know $\emerge(\mathrm{Mary}, \mathrm{red}, d)$
for all $d$?

The answer is subtle.  Mary knows the \emph{axioms} of the ideatic
space and therefore knows that $\emerge(\mathrm{Mary}, \mathrm{red}, d)
= \rs(\mathrm{Mary} \compose \mathrm{red}, d) - \rs(\mathrm{Mary}, d) -
\rs(\mathrm{red}, d)$.  But she does not know the \emph{value} of
this expression until she actually performs the composition---until
she actually sees red.  Knowledge of the axioms is not knowledge of
all their instances.  The axioms constrain the emergence but do not
determine it without the specific compositional act.

This is analogous to knowing the laws of physics without knowing the
initial conditions.  Mary knows the law ($\emerge = \rs \circ \compose -
\rs - \rs$) but not the initial condition (the specific resonance
values).  The ``new knowledge'' she gains upon seeing red is the
specific value of the emergence function at the specific arguments
$(\mathrm{Mary}, \mathrm{red}, d)$.

\subsection{The Inverted Qualia Objection}

Could two ideas have the same emergence profile but ``inverted''
qualia---experiencing red where the other experiences green?  DP28--30
rule this out.  If two ideas have the same emergence profile, then
they contribute identically to all compositions (DP28), differ only
in individual resonance (DP29), and have the same experience of every
encounter (DP30).  There is no room for inversion because there is
no independent ``quale'' to invert.

The inverted qualia scenario depends on the assumption that qualia
are independent of functional role.  DP28--30 deny this assumption:
the emergence profile (functional role) exhaustively characterizes
the experiential profile.  If $\emerge(a, c, d) = \emerge(b, c, d)$
for all $c, d$, then $a$ and $b$ experience the same ``what it is
like'' in every context.  There is nothing left to invert.

\begin{reflectionbox}
\noindent\textbf{Reflection: Why Zombies Seem Conceivable.}

If DP28 proves that zombies are impossible, why do they seem conceivable?
Chalmers's argument relies on the apparent conceivability of a being
that functions like a conscious being but lacks experience.

The answer, from the IDT perspective, is that the conceivability is an
illusion generated by an incomplete understanding of the algebraic
structure.  When we \emph{imagine} a zombie, we implicitly strip away the
emergence profile while keeping the resonance profile.  But DP28--30 show
that you cannot do this consistently.  The emergence profile is not an
independent component that can be added or removed; it is \emph{constitutive}
of the idea's compositional behavior.

Zombies are ``conceivable'' only in the way that a round square is
conceivable: you can form the words, but you cannot consistently
instantiate the concept.  The algebraic structure of the ideatic space
prevents it.
\end{reflectionbox}

\section{Philosophical Coda}

The zombie theorems DP28--30 constitute what is perhaps the most
philosophically provocative result in this volume.  They say, in
mathematical language, that consciousness---understood as the
emergence profile of an idea---is not an optional add-on to
functional organization.  It is functional organization.  There is
nothing over and above the emergence, nothing beyond the
compositional role, nothing hidden behind the functional facade.

This is not eliminativism.  We are not saying that consciousness
does not exist.  We are saying that consciousness \emph{is}
emergence, and emergence is a precisely defined, axiomatically
constrained mathematical quantity.  The hard problem of
consciousness dissolves because the supposed gap between function
and experience was never a real gap---it was an artifact of trying to
describe a single algebraic structure in two different vocabularies.

Chalmers is right that consciousness is something remarkable.
He is wrong that it is something \emph{extra}.  The emergence
function $\emerge(a, c, d) = \rs(a \compose c, d) - \rs(a, d) - \rs(c, d)$
is both the objective, third-person description of how ideas compose
and the subjective, first-person description of what it is like to
be an idea in the presence of others.  These are not two descriptions
of two different things.  They are one description of one thing.

The zombie is dead.  Long live emergence.


%%%%%%%%%%%%%%%%%%%%%%%%%%%%%%%%%%%%%%%%%%%%%%%%%%%%%%%%%%%%%%%%%%%%%%%%
%%  CHAPTER 9 — ASPECT PERCEPTION: WHY THE SAME IDEA LOOKS DIFFERENT
%%%%%%%%%%%%%%%%%%%%%%%%%%%%%%%%%%%%%%%%%%%%%%%%%%%%%%%%%%%%%%%%%%%%%%%%

\chapter{Aspect Perception --- Why the Same Idea Looks Different}
\label{ch:aspect-perception}

\epigraph{``I contemplate a face, and then suddenly notice its likeness to
another.  I see that it has not changed; and yet I see it differently.
I call this experience `noticing an aspect.'\,''}{Ludwig Wittgenstein, \textit{Philosophical Investigations}, Part~II, \S xi}

\section{Wittgenstein's Duck-Rabbit}

In Part~II of the \textit{Philosophical Investigations}, Wittgenstein
introduces what has become one of the most famous images in philosophy:
the duck-rabbit.  The drawing can be seen as a duck (facing left) or
as a rabbit (facing right).  Nothing about the drawing changes when
you switch from one interpretation to the other.  The lines remain the
same, the angles remain the same, the physical stimulus is identical.
Yet the \emph{experience} is radically different.

Wittgenstein calls this ``noticing an aspect'' or ``seeing-as.''  It is
not a matter of interpretation in the theoretical sense---you do not
\emph{infer} that the drawing is a rabbit after previously thinking it
was a duck.  The switch is immediate, gestalt, and involuntary.  One
moment you see a duck; the next moment you see a rabbit.  The question
is: what has changed?

The standard philosophical answer is that the \emph{frame} has changed.
You bring a certain perceptual framework to the drawing, and the
framework determines what you see.  But this answer merely pushes the
question back: what is a frame, and how does it determine perception?

In this chapter, we formalize the structure of aspect perception using
the emergence function and prove three theorems that characterize the
mechanics of the gestalt switch.  The key insight is that the
difference between two aspects of the same observation is \emph{exactly}
the difference in emergence between two frames.  The gestalt switch
is not a mysterious qualitative change but a quantitative difference
in the emergence function.

\section{Frames and Observations}

\begin{definition}[Framed Observation]
\label{def:framed-observation}
Let $f \in \IS$ be a \textbf{frame} (a perceptual or conceptual
framework) and $\mathrm{obs} \in \IS$ be an \textbf{observation}
(a perceptual stimulus, a text, an idea).  The \textbf{framed
observation} is the composition $f \compose \mathrm{obs}$.
\end{definition}

The framed observation is what you \emph{experience} when you encounter
$\mathrm{obs}$ through the lens of frame $f$.  Different frames
produce different experiences of the same observation, just as the
duck-frame and the rabbit-frame produce different experiences of the
same drawing.

The resonance of the framed observation with any idea $c$ is:
\[
\rs(f \compose \mathrm{obs}, c) = \rs(f, c) + \rs(\mathrm{obs}, c) +
\emerge(f, \mathrm{obs}, c).
\]
This decomposition is fundamental.  The total resonance has three
components: the frame's direct resonance with $c$, the observation's
direct resonance with $c$, and the emergence---the genuinely new
content that arises from the encounter of frame and observation.

\section{The Aspect Difference Theorem (DP31)}

\begin{theorem}[Aspect Difference Decomposition --- DP31]
\label{thm:aspect-difference}
For any two frames $f_1, f_2 \in \IS$, any observation
$\mathrm{obs} \in \IS$, and any idea $c \in \IS$:
\[
\rs(f_1 \compose \mathrm{obs}, c) - \rs(f_2 \compose \mathrm{obs}, c) =
\bigl(\rs(f_1, c) - \rs(f_2, c)\bigr) +
\bigl(\emerge(f_1, \mathrm{obs}, c) - \emerge(f_2, \mathrm{obs}, c)\bigr).
\]
The difference between two framings of the same observation decomposes
into exactly two components: the direct frame difference and the
emergence difference.
\end{theorem}

\begin{proof}
By the score decomposition (derived from the axioms):
\begin{align*}
\rs(f_1 \compose \mathrm{obs}, c) &= \rs(f_1, c) + \rs(\mathrm{obs}, c) +
\emerge(f_1, \mathrm{obs}, c), \\
\rs(f_2 \compose \mathrm{obs}, c) &= \rs(f_2, c) + \rs(\mathrm{obs}, c) +
\emerge(f_2, \mathrm{obs}, c).
\end{align*}
Subtracting the second from the first, the $\rs(\mathrm{obs}, c)$ terms
cancel, yielding the result.
\end{proof}

\begin{lstlisting}
-- From IDT_v3b_DeepPhilosophy.lean

/-- DP31: Aspect Difference Decomposition. -/
theorem aspect_difference (f₁ f₂ obs c : I) :
    rs (compose f₁ obs) c - rs (compose f₂ obs) c =
    (rs f₁ c - rs f₂ c) +
    (emergence f₁ obs c - emergence f₂ obs c) := by
  have h1 := rs_compose_eq f₁ obs c
  have h2 := rs_compose_eq f₂ obs c
  linarith
\end{lstlisting}

\begin{surprisebox}
\noindent\textbf{Why this is surprising.}
The observation $\mathrm{obs}$ drops out of the difference entirely.
When you compare two framings of the \emph{same} observation, the
observation's direct resonance cancels.  The only things that matter
are (1) how the frames themselves resonate with $c$, and (2) how
the frames generate different emergences when combined with the
observation.

The ``gestalt switch''---the sudden flip from seeing a duck to
seeing a rabbit---is precisely the second term:
$\emerge(f_1, \mathrm{obs}, c) - \emerge(f_2, \mathrm{obs}, c)$.
This is the emergence difference between the two framings, and it
captures exactly what changes when you ``notice an aspect.''
The lines on the page don't change; the emergence does.
\end{surprisebox}

\begin{interpretation}[Kuhn's Paradigm Shifts]
\label{interp:kuhn-paradigm}
Thomas Kuhn's \textit{The Structure of Scientific Revolutions} (1962)
argues that scientific revolutions are gestalt switches at the
communal level.  When Copernicus replaced Ptolemy, scientists did not
merely accept a new theory; they saw the same observations
\emph{differently}.  The sun's motion across the sky, which had been
``seen as'' the sun orbiting the earth, was now ``seen as'' the
earth rotating.  The observations didn't change; the frame did.

DP31 gives this a precise mathematical formulation.  Let $f_1$ be the
Ptolemaic frame and $f_2$ the Copernican frame.  Let $\mathrm{obs}$ be
the observation of the sun's apparent motion.  The difference between
the two paradigms' assessments of this observation, relative to any
theoretical consequence $c$, is:
\[
\bigl(\rs(f_1, c) - \rs(f_2, c)\bigr) +
\bigl(\emerge(f_1, \mathrm{obs}, c) - \emerge(f_2, \mathrm{obs}, c)\bigr).
\]
The first term captures the direct theoretical difference between the
paradigms.  The second term---the emergence difference---captures the
genuinely new content that each paradigm generates from the observation.
This is the ``incommensurability'' that Kuhn famously (and controversially)
attributed to paradigm shifts: the two paradigms literally produce
different emergences from the same observations.

Kuhn's critics (Popper, Lakatos, Laudan) objected that paradigms could
not be truly incommensurable because scientists can compare predictions.
DP31 resolves this debate: the paradigms \emph{are} commensurable in
the sense that the difference is a well-defined real number.  But they
\emph{are} incommensurable in the sense that the emergence terms can
differ, and emergence cannot be reduced to direct resonance.
\end{interpretation}

\section{Linear Frames Cannot Create Aspects (DP32)}

\begin{definition}[Left-Linear Idea]
\label{def:left-linear}
An idea $f \in \IS$ is \textbf{left-linear} if composing $f$ with
any idea $b$ produces zero emergence:
\[
\emerge(f, b, c) = 0 \quad \text{for all } b, c \in \IS.
\]
Equivalently, $\rs(f \compose b, c) = \rs(f, c) + \rs(b, c)$ for
all $b, c$.
\end{definition}

A left-linear idea is one that contributes to compositions
\emph{additively}---without generating any new content beyond what the
components individually provide.  It is a ``transparent'' frame: it
does not color, distort, or enhance the observation.  It merely adds
its own resonance to the observation's resonance.

\begin{theorem}[Linear Frames Cannot Create Aspects --- DP32]
\label{thm:linear-frames-no-aspects}
If $f_1$ and $f_2$ are both left-linear, then for any observation
$\mathrm{obs}$ and idea $c$:
\[
\rs(f_1 \compose \mathrm{obs}, c) - \rs(f_2 \compose \mathrm{obs}, c) =
\rs(f_1, c) - \rs(f_2, c).
\]
The difference between two linear framings reduces to the direct frame
difference.  There is no emergence difference, and therefore no
gestalt switch.
\end{theorem}

\begin{proof}
If $f_1$ is left-linear, then $\emerge(f_1, \mathrm{obs}, c) = 0$
by definition.  Similarly for $f_2$.  Substituting into DP31:
\[
\rs(f_1 \compose \mathrm{obs}, c) - \rs(f_2 \compose \mathrm{obs}, c) =
(\rs(f_1, c) - \rs(f_2, c)) + (0 - 0) = \rs(f_1, c) - \rs(f_2, c).
\]
\end{proof}

\begin{lstlisting}
-- From IDT_v3b_DeepPhilosophy.lean

/-- DP32: Linear Frames Cannot Create Aspects. -/
theorem linear_frames_no_aspects (f₁ f₂ : I)
    (h1 : leftLinear f₁) (h2 : leftLinear f₂) (obs c : I) :
    rs (compose f₁ obs) c - rs (compose f₂ obs) c =
    rs f₁ c - rs f₂ c := by
  have hf1 := leftLinear_additive f₁ h1 obs c
  have hf2 := leftLinear_additive f₂ h2 obs c
  linarith
\end{lstlisting}

\begin{surprisebox}
\noindent\textbf{Why this is surprising.}
Linear frames cannot produce gestalt switches.  If both frames are
additive---if they contribute to compositions without generating any
emergence---then the difference between the two framings is just
the difference between the frames themselves.  There is no ``extra''
content, no surprise, no aspect perception.

This has a striking consequence for AI: systems based on linear models
(linear regression, linear classifiers, even some neural networks
without sufficient nonlinearity) are \emph{provably incapable} of
aspect perception.  They can distinguish frames (the first term in DP31)
but cannot generate gestalt switches (the second term).

To see the rabbit where you once saw the duck, you need nonlinearity.
You need emergence.  You need a frame that does not merely add its
content to the observation but \emph{transforms} the observation in a
way that generates genuinely new resonance.  Linear thinking is
duck-thinking: it sees only one aspect, forever.
\end{surprisebox}

\begin{interpretation}[Nonlinearity and Creative Insight]
\label{interp:nonlinearity-insight}
The theorem suggests a deep connection between nonlinearity and
creative insight.  Creativity---the ability to see old things in new
ways, to make connections that were not previously visible, to have
``aha moments''---requires nonlinear frames.  A creative mind is one
whose composition with observations generates large emergence:
$\emerge(f_{\mathrm{creative}}, \mathrm{obs}, c) \neq 0$ for many
$c$.

This is consistent with research in cognitive science on the role
of analogical thinking, bisociation (Koestler), and lateral thinking
(de Bono) in creativity.  All of these processes involve bringing
a \emph{surprising} frame to an observation---a frame that generates
emergence precisely because it is not the ``obvious'' frame.

The mathematically precise claim is that linearity kills creativity.
If a thinker's frame is left-linear, then composing the frame with
any observation produces zero emergence, and the thinker can never
experience a gestalt switch.  The condition for aspect perception---and
therefore for creativity---is the existence of nonlinear frames.
\end{interpretation}

\section{The Gestalt Switch Is Bounded (DP33)}

\begin{theorem}[Gestalt Switch Bounded --- DP33]
\label{thm:gestalt-switch-bounded}
For any frames $f_1, f_2$, observation $\mathrm{obs}$, and idea $c$:
\begin{align}
\emerge(f_1, \mathrm{obs}, c)^2 &\leq
\mathrm{weight}(f_1 \compose \mathrm{obs}) \cdot \mathrm{weight}(c), \\
\emerge(f_2, \mathrm{obs}, c)^2 &\leq
\mathrm{weight}(f_2 \compose \mathrm{obs}) \cdot \mathrm{weight}(c).
\end{align}
The magnitude of the emergence from any framing is bounded by the
geometric mean of the framed observation's weight and the observer's weight.
\end{theorem}

\begin{proof}
Both inequalities are direct applications of Axiom~E4
(emergence bounded), with the substitution $a = f_i$, $b = \mathrm{obs}$,
$c = c$:
\[
\emerge(f_i, \mathrm{obs}, c)^2 \leq
\rs(f_i \compose \mathrm{obs}, f_i \compose \mathrm{obs}) \cdot
\rs(c, c) = \mathrm{weight}(f_i \compose \mathrm{obs}) \cdot
\mathrm{weight}(c).
\]
\end{proof}

\begin{lstlisting}
-- From IDT_v3b_DeepPhilosophy.lean

/-- DP33: The Gestalt Switch Is Bounded. -/
theorem gestalt_switch_bounded (f₁ f₂ obs c : I) :
    (emergence f₁ obs c) ^ 2 ≤
      weight (compose f₁ obs) * weight c ∧
    (emergence f₂ obs c) ^ 2 ≤
      weight (compose f₂ obs) * weight c := by
  unfold weight
  exact ⟨emergence_bounded' f₁ obs c, emergence_bounded' f₂ obs c⟩
\end{lstlisting}

\begin{surprisebox}
\noindent\textbf{Why this is surprising.}
Even the most dramatic gestalt switch has a mathematical ceiling.
The emergence from any framing cannot exceed
$\sqrt{\mathrm{weight}(f \compose \mathrm{obs}) \cdot \mathrm{weight}(c)}$.
This means that the ``aha!'' of seeing the rabbit cannot be
infinitely surprising---it is bounded by the weight of the
framed observation and the weight of the observer.

This has implications for the phenomenology of surprise.  We know
from experience that some gestalt switches are more dramatic than
others: seeing the hidden face in a visual puzzle is more surprising
than noticing that a cloud looks like a sheep.  DP33 says that this
variation is lawful: more dramatic switches require either a heavier
framed observation (more context) or a heavier observer (more
background knowledge).

The bound also implies that lightweight ideas---ideas with low
self-resonance---cannot participate in dramatic gestalt switches.
Surprise requires substance.
\end{surprisebox}

\begin{interpretation}[The Phenomenology of Surprise]
\label{interp:phenomenology-surprise}
Husserl's phenomenology distinguishes between ``empty'' and ``filled''
intentions.  An empty intention is a thought directed at an object
that is not present (thinking about the Eiffel Tower while sitting
in New York).  A filled intention is a thought that encounters its
object directly (seeing the Eiffel Tower in Paris).  The transition
from empty to filled intention---the moment of ``fulfillment''---is
accompanied by a distinctive phenomenal quality: the ``aha!'' of
recognition.

DP33 puts a bound on this ``aha!'' moment.  The emergence from
fulfillment (composing the expectation-frame with the actual
observation) cannot exceed the square root of the product of the
weights.  Lightweight expectations produce lightweight fulfillments.
Only a rich, weighty anticipation can generate a rich, weighty
``aha!''

This is consistent with the everyday observation that surprise requires
expectation.  You cannot be surprised if you had no prior frame.
And the depth of the surprise is bounded by the depth of the frame
and the depth of the observation.

The mathematics of surprise, it turns out, is a Cauchy--Schwarz
inequality.  The inner product structure of resonance constrains
the magnitude of emergence, and emergence \emph{is} surprise.
\end{interpretation}

\section{Aspect Perception and the Structure of Understanding}

The three theorems of this chapter (DP31--33) constitute a complete
mathematical theory of aspect perception.

DP31 tells us what the gestalt switch \emph{is}: the emergence
difference between two frames applied to the same observation.
This is a precise, quantitative characterization of a phenomenon
that philosophers have struggled to describe for over a century.

DP32 tells us what is \emph{necessary} for the gestalt switch:
nonlinearity.  Linear frames produce no emergence and therefore no
aspect perception.  The capacity for gestalt switches---and therefore
for creativity, metaphor, and insight---depends on the existence of
nonlinear frames in the ideatic space.

DP33 tells us what \emph{constrains} the gestalt switch: the
Cauchy--Schwarz bound.  The magnitude of surprise is limited by the
weight of the composition and the weight of the observer.  Dramatic
insights require substantial context.

Together, these theorems explain why the duck-rabbit is possible,
why it is bounded, and why it requires a certain kind of mind
(a nonlinear one) to experience it.

\begin{reflectionbox}
\noindent\textbf{Reflection: Seeing-As and Being-In.}

Wittgenstein's analysis of aspect perception is often treated as a
marginal topic in his philosophy---a curiosity about visual
perception that has little bearing on the deeper issues of language,
meaning, and form of life.  The theorems of this chapter suggest
otherwise.

Aspect perception is not a special case; it is the \emph{general}
case.  Every act of understanding involves a frame, an observation,
and emergence.  Every understanding is a ``seeing-as''---a framing
of the world through a particular lens that generates particular
emergences.  The duck-rabbit is merely the most vivid illustration
of a universal structure.

To understand is always to see \emph{as}.  And the mathematics of
emergence tells us that this ``as'' is not arbitrary but lawful,
bounded, and structurally constrained.
\end{reflectionbox}

\section{Philosophical Coda}

The results of this chapter illuminate a deep connection between
three apparently unrelated phenomena: Wittgenstein's aspect perception,
Kuhn's paradigm shifts, and the role of nonlinearity in creative
thinking.

All three are instances of the same mathematical structure: a change
of frame that produces a change in emergence.  The duck-rabbit switch
is a local gestalt change; the paradigm shift is a global gestalt
change; creative insight is a personal gestalt change.  In each case,
the mechanism is the same: a new frame generates new emergence from
the same observation, and the emergence difference is the experience
of seeing something new.

The linearity condition (DP32) reveals why some minds---and some
computational architectures---are incapable of these switches.
A purely linear system, no matter how sophisticated, cannot generate
emergence.  It can process information, combine features, and compute
outputs, but it cannot experience the ``aha!'' of seeing the rabbit
in the duck.

This is not a limitation of hardware or software; it is a
mathematical constraint.  The capacity for aspect perception, creative
insight, and paradigm shifts is the capacity for nonlinear composition,
and nonlinear composition is the source of emergence.

The deepest conclusion of this chapter is that emergence---the
non-additive residual of composition---is not an incidental feature
of the ideatic space but its \emph{central drama}.  Emergence is what
makes the gestalt switch possible, what makes paradigm shifts
transformative, and what makes creativity creative.  Without
emergence, composition would be merely additive, frames would be
merely transparent, and understanding would be merely aggregation.

With emergence, everything changes.  And the change is bounded.


%%%%%%%%%%%%%%%%%%%%%%%%%%%%%%%%%%%%%%%%%%%%%%%%%%%%%%%%%%%%%%%%%%%%%%%%
%%  CHAPTER 10 — THE RULE-FOLLOWING PARADOX: EVERY RULE HAS INFINITE
%%                INTERPRETATIONS
%%%%%%%%%%%%%%%%%%%%%%%%%%%%%%%%%%%%%%%%%%%%%%%%%%%%%%%%%%%%%%%%%%%%%%%%

\chapter{The Rule-Following Paradox --- Every Rule Has Infinite Interpretations}
\label{ch:rule-following}

\epigraph{``This was our paradox: no course of action could be
determined by a rule, because every course of action can be made out to
accord with the rule.  The answer was: if everything can be made out to
accord with the rule, then it can also be made out to conflict with it.
And so there would be neither accord nor conflict here.''}{Ludwig Wittgenstein, \textit{Philosophical Investigations}, \S 201}

\section{Wittgenstein and the Paradox of Rules}

In one of the most celebrated and contentious passages in twentieth-century
philosophy, Wittgenstein presents the rule-following paradox.  Consider
the simplest possible rule: ``add 2.''  You have been computing the
sequence $2, 4, 6, 8, 10, \ldots$ and now you reach $1000$.  The obvious
next step is $1002$.  But consider an alternative rule: ``quus,'' defined
as follows:
\[
x \oplus y = \begin{cases}
x + y & \text{if } x, y < 1000, \\
5 & \text{otherwise.}
\end{cases}
\]
Your entire past behavior is consistent with both ``plus'' and ``quus.''
Nothing in your previous computations---not even your mental states---can
distinguish which rule you were following.  And yet ``plus'' gives $1002$
while ``quus'' gives $5$.

This is Saul Kripke's reconstruction of Wittgenstein's paradox, from
\textit{Wittgenstein on Rules and Private Language} (1982).  Kripke
argues that there is ``no fact of the matter'' about which rule you are
following.  Every finite sequence of actions is compatible with infinitely
many rules, and no mental state can fix which one is operative.

The paradox threatens to undermine the very concept of meaning.  If no
fact determines which rule you are following, then there is no fact about
what your words mean, no fact about what you intend, and no fact about
whether you are computing correctly.

In this chapter, we show that the IDT axioms illuminate the structure
of the paradox.  Rule application always involves emergence, and
emergence is the ``interpretation'' that the rule itself does not
determine.  The communitarian solution---the idea that rule-following
is grounded in shared practice---receives a precise mathematical
formulation.

\section{Rule Application as Composition}

\begin{definition}[Rule Application]
\label{def:rule-application}
A \textbf{rule} is an idea $r \in \IS$.  \textbf{Applying} rule $r$
to input $a$ is composition:
\[
\mathrm{apply}(r, a) := r \compose a.
\]
\end{definition}

This identification is natural.  A rule is a procedure that takes an
input and produces an output.  In the ideatic space, the only way
to combine ideas is composition, so applying a rule to an input is
composing the rule-idea with the input-idea.

The resonance of the output $r \compose a$ with any observer $c$
decomposes as:
\[
\rs(r \compose a, c) = \rs(r, c) + \rs(a, c) + \emerge(r, a, c).
\]
The output resonance has three sources: the rule's inherent resonance,
the input's inherent resonance, and the emergence---the non-additive
residual that arises from the specific encounter of rule and input.

\section{Rule Application Creates Emergence (DP34)}

\begin{theorem}[Rule Creates Emergence --- DP34]
\label{thm:rule-creates-emergence}
For any rule $r$, input $a$, and observer $c$:
\begin{align}
\rs(r \compose a, c) &= \rs(r, c) + \rs(a, c) + \emerge(r, a, c), \\
\emerge(r, a, c) &= \rs(r \compose a, c) - \rs(r, c) - \rs(a, c).
\end{align}
The rule does not fully determine the output.  There is always an
emergence term that depends on the specific input and observer.
\end{theorem}

\begin{proof}
The first equation is the score decomposition, which follows from the
definition of emergence:
$\emerge(r, a, c) = \rs(r \compose a, c) - \rs(r, c) - \rs(a, c)$.
Rearranging gives $\rs(r \compose a, c) = \rs(r, c) + \rs(a, c) + \emerge(r, a, c)$.
The second equation is the definition of emergence itself.
\end{proof}

\begin{lstlisting}
-- From IDT_v3b_DeepPhilosophy.lean

/-- DP34: Rule Application Creates Emergence. -/
theorem rule_creates_emergence (r a c : I) :
    rs (compose r a) c = rs r c + rs a c + emergence r a c ∧
    emergence r a c = rs (compose r a) c - rs r c - rs a c := by
  constructor
  · exact rs_compose_eq r a c
  · unfold emergence; ring
\end{lstlisting}

\begin{surprisebox}
\noindent\textbf{Why this is surprising.}
The rule $r$ does not fully determine the output of its application.
Even if you know $\rs(r, c)$ for all $c$---the complete resonance
profile of the rule---you cannot predict $\rs(r \compose a, c)$
without knowing the emergence $\emerge(r, a, c)$.  The emergence
depends on $r$, $a$, \emph{and} $c$---the rule, the input, and the
observer.

This is the mathematical content of Wittgenstein's paradox.  The rule
alone does not determine its application because the application always
involves emergence, and emergence is not derivable from the rule in
isolation.  Every rule application is an \emph{act of interpretation}
(an emergence event), not a mechanical deduction.

The ``quus'' paradox is demystified: ``plus'' and ``quus'' can agree
on all past inputs because the emergence term absorbs the difference.
The rule's resonance profile is the same for both (by hypothesis), but
the emergence from applying the rule to input $1000$ differs.
\end{surprisebox}

\begin{interpretation}[The Rule as Under-Determined]
\label{interp:rule-underdetermined}
Kripke's Wittgenstein asks: what fact about me determines that I mean
\emph{plus} rather than \emph{quus}?  The standard answers---my
mental state, my dispositions, my past behavior---all fail, because
they are all compatible with both rules.

DP34 provides a new answer: no \emph{single} fact determines which rule
I follow.  The output of rule application is determined by three
components: the rule's resonance ($\rs(r, c)$), the input's resonance
($\rs(a, c)$), and the emergence ($\emerge(r, a, c)$).  The rule alone
fixes only the first component.  The other two depend on the context
of application.

This is not skepticism about rules.  It is a positive thesis about
the structure of rule-following: rule-following is always
\emph{contextual}.  The rule provides a framework, but the framework
must be instantiated in a specific context (input, observer), and the
instantiation generates emergence.

In Wittgenstein's terms: the rule does not ``contain'' its applications.
Each application is a new composition, and each composition generates
new emergence.  This is not a defect of rules but a feature of
the ideatic space.
\end{interpretation}

\section{Community Agreement Requires Emergence (DP35)}

\begin{theorem}[Community Requires Emergence --- DP35]
\label{thm:community-requires-emergence}
For any two rules $r_1, r_2$, input $a$, and observer $c$:
\[
\rs(r_1 \compose a, c) = \rs(r_2 \compose a, c) \iff
\rs(r_1, c) + \emerge(r_1, a, c) = \rs(r_2, c) + \emerge(r_2, a, c).
\]
Two agents agree on the output of rule application if and only if
their individual resonances plus emergences are aligned.
\end{theorem}

\begin{proof}
($\Rightarrow$): If $\rs(r_1 \compose a, c) = \rs(r_2 \compose a, c)$,
then by the score decomposition:
\[
\rs(r_1, c) + \rs(a, c) + \emerge(r_1, a, c) =
\rs(r_2, c) + \rs(a, c) + \emerge(r_2, a, c).
\]
Canceling $\rs(a, c)$ gives
$\rs(r_1, c) + \emerge(r_1, a, c) = \rs(r_2, c) + \emerge(r_2, a, c)$.

($\Leftarrow$): If $\rs(r_1, c) + \emerge(r_1, a, c) = \rs(r_2, c) +
\emerge(r_2, a, c)$, then adding $\rs(a, c)$ to both sides and applying
the score decomposition gives
$\rs(r_1 \compose a, c) = \rs(r_2 \compose a, c)$.
\end{proof}

\begin{lstlisting}
-- From IDT_v3b_DeepPhilosophy.lean

/-- DP35: Community Agreement Requires Shared Emergence. -/
theorem community_requires_emergence (r₁ r₂ a c : I) :
    rs (compose r₁ a) c = rs (compose r₂ a) c ↔
    rs r₁ c + emergence r₁ a c =
    rs r₂ c + emergence r₂ a c := by
  constructor
  · intro h
    have h1 := rs_compose_eq r₁ a c
    have h2 := rs_compose_eq r₂ a c
    linarith
  · intro h
    have h1 := rs_compose_eq r₁ a c
    have h2 := rs_compose_eq r₂ a c
    linarith
\end{lstlisting}

\begin{surprisebox}
\noindent\textbf{Why this is surprising.}
Two agents can follow ``different rules'' ($r_1 \neq r_2$) and yet
agree on the output for a given input.  This happens when the
difference in their rule resonances is exactly compensated by the
difference in their emergences:
\[
\rs(r_1, c) - \rs(r_2, c) = \emerge(r_2, a, c) - \emerge(r_1, a, c).
\]
Conversely, two agents can follow the ``same rule'' ($r_1 = r_2$) and
yet disagree on the output, if their emergence profiles differ.

This theorem reveals that ``following the same rule'' is not a
syntactic notion (having the same rule-idea) but a \emph{semantic}
notion (producing the same output in context).  Two people can mean
different things by ``plus'' and yet agree on all computations, as
long as their emergences compensate for their rule differences.

This is exactly the communitarian solution to the rule-following
paradox: what makes us follow the ``same rule'' is not that we
have identical mental representations but that we produce the
same outputs in the same contexts.  And producing the same outputs
requires alignment of resonance \emph{and} emergence.
\end{surprisebox}

\begin{interpretation}[Kripke's Communitarian Solution]
\label{interp:kripke-communitarian}
Kripke concludes his book with what he calls a ``skeptical solution''
to the rule-following paradox: there is no fact about what rule an
individual follows, but there are facts about what a \emph{community}
agrees on.  Rule-following is grounded in communal agreement, not
individual mental states.

DP35 gives this solution a precise mathematical form.  Community
agreement on rule application requires:
\[
\rs(r_1, c) + \emerge(r_1, a, c) = \rs(r_2, c) + \emerge(r_2, a, c).
\]
This is a condition on the \emph{sum} of resonance and emergence, not
on either alone.  Two agents with different rules and different
emergences can agree, and two agents with the same rule can disagree.

The communitarian solution works because what matters for agreement
is the compositional output $\rs(r \compose a, c)$, not the
decomposition into $\rs(r, c)$ and $\emerge(r, a, c)$.  The
community doesn't check your internal rule; it checks your external
output.  And the output is determined by the \emph{sum} of resonance
and emergence, which is a holistic, context-dependent quantity.

This suggests that Kripke's skepticism is too strong.  There \emph{is}
a fact about what rule you follow---it is the fact about which
$(r, a) \mapsto \rs(r \compose a, c)$ function you instantiate.  But
this fact is not reducible to the rule alone; it involves the
emergence, which depends on the context of application.
\end{interpretation}

\section{Rule Composition and Higher-Order Emergence (DP36)}

\begin{theorem}[Rule Composition Creates Higher-Order Emergence --- DP36]
\label{thm:rule-composition-higher-order}
For any rules $r_1, r_2$, input $a$, and observer $d$:
\[
\emerge(r_1, r_2, d) + \emerge(r_1 \compose r_2, a, d) =
\emerge(r_2, a, d) + \emerge(r_1, r_2 \compose a, d).
\]
The emergences from composing two rules and then applying the composition
satisfy a cocycle condition.  Rules interact non-additively.
\end{theorem}

\begin{proof}
This is the cocycle condition (derived from the axioms) applied to
$a = r_1$, $b = r_2$, $c = a$, $d = d$:
\[
\emerge(r_1, r_2, d) + \emerge(r_1 \compose r_2, a, d) =
\emerge(r_2, a, d) + \emerge(r_1, r_2 \compose a, d).
\]
\end{proof}

\begin{lstlisting}
-- From IDT_v3b_DeepPhilosophy.lean

/-- DP36: Rule Composition Creates Higher-Order Emergence. -/
theorem rule_composition_higher_order (r₁ r₂ a d : I) :
    emergence r₁ r₂ d + emergence (compose r₁ r₂) a d =
    emergence r₂ a d + emergence r₁ (compose r₂ a) d := by
  exact cocycle_condition r₁ r₂ a d
\end{lstlisting}

\begin{surprisebox}
\noindent\textbf{Why this is surprising.}
When you compose two rules and then apply the composition to an input,
the resulting emergence is \emph{not} the sum of the individual
emergences.  There is a correction term---the cocycle---that accounts
for the interaction between the rules.

Consider what this means concretely.  Suppose $r_1$ is ``take the
square root'' and $r_2$ is ``add one.''  The composition $r_1 \compose r_2$
is ``add one, then take the square root.''  When you apply this
composite rule to an input $a$, the emergence is NOT the emergence from
adding one plus the emergence from taking the square root.  The two
operations interact: the emergence from taking the square root of
$a + 1$ is different from the emergence from taking the square root
of $a$, and the difference is precisely the cocycle correction.

This is the mathematical expression of a deep truth about rules:
\emph{rules do not compose simply}.  Every composition of rules
introduces a higher-order interaction that cannot be predicted from
the individual rules alone.  Rule-following is irreducibly contextual,
and the context includes not just the input but the other rules in
the composition.
\end{surprisebox}

\begin{interpretation}[The Practice Behind the Rule]
\label{interp:practice-behind-rule}
Wittgenstein's solution to the rule-following paradox (in the later
sections of the \textit{Philosophical Investigations}) is often
summarized as: understanding a rule is mastering a \emph{practice}.
The rule does not contain its own application; it must be embedded
in a form of life, a network of practices, customs, and
institutions that give it meaning.

DP36 formalizes this insight.  The cocycle condition shows that the
practice behind a rule is not a single emergence but a structured
system of emergences that satisfy an algebraic constraint.  When
two rules are composed, the emergences interact according to the
cocycle identity.  This interaction is the ``practice''---the
shared background of compositional experience that makes rule-following
possible.

The cocycle condition has a cohomological interpretation (developed
in Volume~I): the emergence function is a $2$-cocycle on the
monoid $(\IS, \compose, \void)$ with values in $\R$.  This means
that the ``practice'' of rule-following has the algebraic structure
of a cohomology class---a global constraint on local interactions
that cannot be decomposed into purely local data.

In Wittgenstein's terms: the practice is \emph{holistic}.  You cannot
understand a single rule application in isolation; you must understand
it as part of a system of rule applications, and the system satisfies
the cocycle condition.  The cocycle is the mathematical form of what
Wittgenstein calls a ``form of life.''
\end{interpretation}

\section{The Impossibility of Private Rules}

Combining DP34--36, we can derive a striking consequence: there are
no genuinely ``private'' rules.

A private rule would be a rule $r$ whose application is fully
determined by $r$ alone, without reference to any context or
community.  DP34 shows that this is impossible: the output of rule
application always involves emergence, and emergence depends on the
input and the observer, not just the rule.  There is always an
``interpretation'' step that goes beyond the rule itself.

DP35 shows that community agreement on rule application requires
alignment of both resonance and emergence.  This means that
rule-following is inherently \emph{social}: it requires a shared
context of emergence, not just a shared rule-representation.  Two
individuals can ``share a rule'' only in the sense that their
resonance-plus-emergence profiles agree on the relevant inputs.

DP36 shows that even the composition of rules is not simple but
involves higher-order interactions.  The practice of rule-following
is not a linear sequence of applications but a web of interlocking
emergences that satisfy the cocycle condition.

Together, these results vindicate Wittgenstein's famous dictum: ``To
follow a rule, to make a report, to give an order, to play a game of
chess, are \emph{customs} (uses, institutions)''
(\textit{Philosophical Investigations}, \S 199).  Rules are customs,
and customs are shared emergence profiles.

\begin{reflectionbox}
\noindent\textbf{Reflection: The Mathematics of Custom.}

Wittgenstein's appeal to ``custom'' and ``practice'' has often been
criticized as philosophically unsatisfying---a retreat from explanation
into description.  The theorems of this chapter show that the appeal
to custom is not a retreat but a \emph{discovery}.

Custom has a precise mathematical structure: it is a system of
emergence values that satisfies the cocycle condition.  The cocycle
condition is not a vague hand-wave about ``shared practices''; it is
an exact algebraic constraint that determines how rule applications
interact.

The mathematics of custom is, it turns out, cohomology.  And
cohomology is one of the deepest and most powerful tools in modern
mathematics.  When Wittgenstein appealed to ``forms of life'' as the
ground of rule-following, he was---without knowing it---appealing to
a cohomological structure.

This is perhaps the most remarkable convergence in this entire volume:
the philosopher who distrusted formalization above all others was
groping toward one of the most abstract formalizations in mathematics.
\end{reflectionbox}

\section{Philosophical Coda}

The rule-following paradox has been called the central problem of
the \textit{Philosophical Investigations}---and, by extension, of
all post-Wittgensteinian philosophy.  If no fact determines which
rule you are following, then the foundations of meaning, intention,
and normativity are in jeopardy.

DP34--36 do not resolve the paradox in the sense of providing a
fact that determines the rule.  Instead, they \emph{characterize}
the paradox: rule application always involves emergence, and
emergence is irreducibly contextual.  The rule underdetermines its
application (DP34), community agreement requires shared emergence
(DP35), and rule composition generates higher-order interactions
(DP36).

The solution---if it deserves the name---is that rule-following is
not a matter of mechanically applying a fixed procedure but a matter
of participating in a shared practice whose structure is captured by
the cocycle condition.  The ``fact'' that determines which rule you
are following is not a single fact but an entire \emph{system} of
facts: your emergence profile in all compositional contexts.

This is why rule-following seems paradoxical when viewed through the
lens of classical logic (which seeks single, determinate facts) but
natural when viewed through the lens of IDT (which recognizes that
meaning is always compositional and emergent).

Wittgenstein's paradox is not a paradox.  It is a theorem about the
structure of emergence.  And the theorem says: every rule has infinite
interpretations because every composition generates emergence, and
emergence is the interpretation that the rule itself cannot provide.

The rule-following paradox is, in the end, just the cocycle condition
in philosophical dress.  And the cocycle condition is not a paradox.
It is the law of ideatic composition.
